% =====================================================================
% ESPAÇAMENTOS
% Manual de Normalização do IFPI (2024), seções 3.3, 3.4 e 7.4
%
% \parindent = recuo da primeira linha de cada parágrafo. A norma pede
%              1,25 cm a partir da margem esquerda (seção 3.4, alínea c).
% \parskip   = espaço extra entre parágrafos. A norma pede 0 pt: os
%              parágrafos são separados apenas pelo recuo.
%
% Não ajuste espaçamento parágrafo a parágrafo no texto: mude aqui e o
% documento inteiro acompanha.
% =====================================================================
\setlength{\parindent}{1.25cm}
\setlength{\parskip}{0pt}

% ---------------------------------------------------------------------
% ESPAÇAMENTO ENTRE LINHAS
%
% A regra geral (seção 3.3, alínea a) é 1,5 entre linhas. MAS o artigo
% científico é exceção: "fonte Arial ou Times New Roman, tamanho 12,
% espaçamento simples em todo o artigo" (seções 3.6.1 Nota 02 e 7.4).
%
% Este modelo é um ARTIGO, então o padrão aqui é o espaçamento simples.
% Se você adaptar o modelo para monografia, dissertação ou tese, troque
% \SingleSpacing por \OnehalfSpacing na linha abaixo.
%
% Em qualquer das duas modalidades continuam em espaço simples, por
% conta própria, as citações longas, as notas de rodapé, as referências,
% as legendas e a natureza do trabalho (seção 3.3, alínea b) -- isso o
% abnTeX2 já faz sozinho.
% ---------------------------------------------------------------------
\SingleSpacing
